\section{La trayectoria de un gerente de producto: de los datos estructurados a la generación no estructurada}

Esta sección convierte primero la experiencia del entrevistado en pistas metodológicas. La trayectoria profesional de Wang Guan (王冠) no es solo una presentación de antecedentes: explica por qué pasó de preguntarse «cómo ofrecer capacidades algorítmicas» a preguntarse «cómo diseñar las capacidades de un modelo y los sistemas generativos».

Esta sección examina primero la trayectoria profesional de Wang Guan. Abarca big data, CV y aprendizaje profundo, modelos preentrenados, el puesto de gerente de producto de modelos en Moonshot y la fundación de ONE2X. El valor didáctico de esta trayectoria es que ofrece una historia técnica de la IA desde la perspectiva de un gerente de producto: la IA no es solo la evolución de los modelos algorítmicos, sino también la evolución de la forma de representar los datos, del lugar que ocupa el producto y de la manera en que los usuarios perciben la tecnología.

\begin{figure}[H]
\centering
\includegraphics[width=0.96\textwidth]{product-manager-trajectory.png}
\caption{La trayectoria de Wang Guan en productos de IA: big data, CV y aprendizaje profundo, preentrenamiento, Moonshot y ONE2X. Diagrama conceptual propio, elaborado a partir de la conversación entre 00:02:00 y 00:28:39.}
\end{figure}

\begin{knowledgebox}{Lectura de la figura: el hilo común de la trayectoria es «reducir la barrera de entrada para usar capacidades algorítmicas»}
Desde los perfiles de usuario en Baidu, pasando por la plataforma abierta de algoritmos, los frameworks de código abierto y las herramientas de productividad algorítmica, hasta los productos de modelos en Moonshot, el trabajo de Wang Guan siempre ha girado en torno a «permitir que los usuarios accedan a cierta capacidad algorítmica con una barrera de entrada menor». ONE2X es solo la extensión de esa línea hacia la generación de video y los sistemas generativos.
\end{knowledgebox}

\subsection{Las tres eras de producto de la IA}

Esta subsección extrae de esa experiencia un primer juicio: lo decisivo en el cambio de los ciclos de la IA no es solo el tamaño de los modelos, sino qué forma de datos pueden ajustar.

Wang Guan no divide los ciclos de la IA en 1.0, 2.0 y 3.0, sino según si los datos son estructurados o no. El aprendizaje automático tradicional y la CV temprana dependían de etiquetas estructuradas: perfiles de usuario, cuadros delimitadores, etiquetas de categoría y campos explícitos. Los modelos grandes, en cambio, ajustan datos no estructurados: lenguaje, imágenes, video y representaciones de un mundo continuo. Los datos no estructurados tienen mayor capacidad expresiva y se acercan más a la continuidad del mundo real.

\begin{knowledgebox}{Términos clave: datos estructurados y no estructurados}
\begin{tabularx}{\textwidth}{p{0.22\textwidth}p{0.34\textwidth}X}
\toprule
Término & Explicación breve & Relación con este episodio \\
\midrule
Datos estructurados & Datos con campos, etiquetas, coordenadas o categorías explícitos & Entrada central en la era del aprendizaje automático tradicional y de la CV. \\
Datos no estructurados & Representaciones continuas como texto, imágenes y video, que no se comprimen de antemano en campos fijos & Objeto central de entrenamiento de los modelos grandes y de los sistemas generativos. \\
User profile & Perfil de usuario: describe las preferencias y el comportamiento del usuario mediante etiquetas & Base del modelado del lado del usuario en los sistemas de recomendación. \\
Gerente de producto de modelos & Diseña el desempeño, las capacidades y la evaluación del modelo para que el usuario perciba esas capacidades & Puesto clave que se formó durante la experiencia en Moonshot. \\
\bottomrule
\end{tabularx}
\end{knowledgebox}

\subsection{Por qué los modelos no han reemplazado al gerente de producto}

Wang Guan sostiene que en la era de la IA el valor del gerente de producto ha aumentado, no disminuido. Hay dos razones. Primero, el modelo se parece al System1 humano: comprime de antemano una gran cantidad de información dentro del sistema y puede reaccionar a la entrada de forma rápida e instintiva. El gerente de producto debe diseñar qué capacidades debe tener ese System1, y lograr que esas capacidades se entrenen y que el usuario las perciba. Segundo, la salida del modelo depende de qué tokens la preceden; el gerente de producto debe diseñar el workflow, el framework del agent, la base de conocimiento del dominio y el context, para que el modelo obtenga más tokens útiles al generar.

\begin{importantbox}{Las dos nuevas tareas del gerente de producto de IA}
Primero, definir las capacidades del modelo: qué desempeño vale la pena entrenar, cómo evaluarlo y cómo lo percibe el usuario. Segundo, diseñar el context: cómo convertir el know-how del negocio, los flujos de trabajo, las herramientas y el conocimiento del dominio en tokens útiles que el modelo pueda aprovechar.
\end{importantbox}

\subsection{Resumen de la sección}

La trayectoria de Wang Guan muestra que el gerente de producto de IA no queda absorbido por las capacidades del modelo, sino que pasa a un primer plano. Cuanto más capaz es el modelo, más se necesita a alguien que defina sus capacidades, organice el context, diseñe la forma de las tareas y convierta esas capacidades en un producto que los usuarios puedan entender, por el que estén dispuestos a pagar y que puedan usar de forma continua.
